\section{Solution concepts}
\label{sec:solution-concepts}

This section catalogues the solution concepts the library
formalizes.  All of them sit on top of \TT{KernelGame} or
\TT{GameForm}; concrete language layers (\S\ref{sec:languages})
inherit them via $\toKG$.  The presentation is grouped by the
\emph{deviation family} the concept quantifies over, which is the
underlying organizing principle of the library and makes
several otherwise-disparate equivalences (correlated $\Rightarrow$
coarse correlated, dominance $\Rightarrow$ Nash, etc.) visible as
instances of one schema.

\subsection{The deviation-family schema}
\label{sec:solution-concepts:family}

Fix a game form $\GF$, a player $i$, and a status-quo profile
distribution $\mu \in \PMF{\prod_j \Strat_j}$.  Let
$\Pref$ be a per-player weak preference on $\PMF{\Outcome}$ and
let $\mathcal{D}_i = (D_i, \mathrm{deviate}_i)$ be a
\emph{deviation family} for $i$:\ a type $D_i$ of deviations and a
function $\mathrm{deviate}_i : \PMF{\prod_j \Strat_j} \to D_i \to
\PMF{\prod_j \Strat_j}$.  A profile distribution $\mu$ is
\emph{$\mathcal{D}$-equilibrium for $\Pref$} if no player has a
profitable deviation in $\mathcal{D}$:
\begin{equation}
\label{eq:devschema}
  \forall i,\, \forall d \in D_i.\quad
    \Pref_i\;\bigl(\GF.\mathrm{correlatedOutcome}(\mu),\;
                   \GF.\mathrm{correlatedOutcome}(\mathrm{deviate}_i(\mu, d))\bigr).
\end{equation}
The library identifies exactly two deviation families that matter
for the standard solution-concept inventory:

\begin{itemize}\itemsep2pt
\item $\mathcal{D}^{\mathrm{const}}$:\ $D_i = \Strat_i$, with
$\mathrm{deviate}_i(\mu, s')$ replacing $i$'s coordinate by $s'$
on every sample of $\mu$.

\item $\mathcal{D}^{\mathrm{dep}}$:\ $D_i = \Strat_i \to \Strat_i$,
with $\mathrm{deviate}_i(\mu, \mathrm{dev})$ replacing $i$'s
coordinate $s$ by $\mathrm{dev}(s)$ on every sample.
\end{itemize}

The four ``profile-distribution'' equilibria of the library are
then schema instances:
\[
  \begin{array}{lcl}
    \IsNashFor(\Profile)
      &\equiv& \mathcal{D}^{\mathrm{const}}\text{-eq for } \Dirac_\Profile, \\
    \IsCoarseCorrelatedEqFor(\mu)
      &\equiv& \mathcal{D}^{\mathrm{const}}\text{-eq for } \mu, \\
    \IsCorrelatedEqFor(\mu)
      &\equiv& \mathcal{D}^{\mathrm{dep}}\text{-eq for } \mu.
  \end{array}
\]
Several library-wide consequences follow:
\begin{itemize}
\item \emph{Hierarchy}.  $\mathcal{D}^{\mathrm{const}}$ embeds into
$\mathcal{D}^{\mathrm{dep}}$ via constant deviation functions, so
every $\Pref$-CE is a $\Pref$-CCE.  Setting $\mu = \Dirac_\Profile$
recovers $\Pref$-Nash from $\Pref$-CCE.  Thus pure Nash point masses
are CCEs, and CE distributions are CCE distributions; CE and CCE are
both predicates on profile distributions, while Nash is a predicate
on point profiles.
\item \emph{Monotonicity in $\Pref$}.  If $\Pref^{(2)}$ implies
$\Pref^{(1)}$ pointwise, then any $\mathcal{D}$-equilibrium for
$\Pref^{(2)}$ is a $\mathcal{D}$-equilibrium for $\Pref^{(1)}$.
The library exposes this as
\TT{IsNashFor.mono} and \TT{IsDominantFor.mono}.
\end{itemize}

\leancorr{Schema in \TT{Concepts/Deviation.lean}; full Lean text in
supplement~§3.1.}

\subsection{Nash equilibrium and best response}
\label{sec:solution-concepts:nash}

The EU-specific Nash predicate is
\begin{equation}
\label{eq:nash}
  \IsNash_\KG(\Profile) \;\Longleftrightarrow\;
    \forall i,\, \forall s' \in \Strat_i,\;
    \EU(\KG, \Profile, i) \;\ge\; \EU(\KG, \Profile[i \mapsto s'], i),
\end{equation}
and the preference-parameterized variant on $\GF$ is the schema
instance with $\mu = \Dirac_\Profile$ and $\mathcal{D} =
\mathcal{D}^{\mathrm{const}}$.  A \emph{best response} is a
strategy that maximizes the expected utility against fixed
opponents:
\begin{equation}
\label{eq:br}
  \IsBR_\KG(i, \Profile, s) \;\Longleftrightarrow\;
    \forall s' \in \Strat_i,\;
    \EU(\KG, \Profile[i \mapsto s], i) \;\ge\;
    \EU(\KG, \Profile[i \mapsto s'], i),
\end{equation}
and the standard equivalence
\[
  \IsNash_\KG(\Profile) \;\Longleftrightarrow\;
    \forall i,\, \IsBR_\KG(i, \Profile, \Profile_i)
\]
holds at the $\GF.\IsNashFor$/$\IsBRfor$ level for any
preference and is propagated to the EU level by the bridge of
\S\ref{sec:semantic-core:for-style}.

A \emph{strict Nash equilibrium} requires the inequality to be
strict for every $s' \neq \Profile_i$.  The library's
\TT{StrictNashProperties.lean} records the implication strict
$\Rightarrow$ Nash, asymmetry under a strict preference that is
incompatible with reverse weak preference, and uniqueness in
ordinal-potential games.

\subsection{Dominance}
\label{sec:solution-concepts:dominance}

Two notions of dominance are formalized.  \emph{Weak dominance}:
$s$ weakly dominates $t$ for player $i$ when, for every opponent
profile, $i$'s payoff under $s$ is $\Pref$-preferred to that under
$t$.  \emph{Strict dominance} replaces $\Pref$ by a strict
preference $\sPref$ (typically the strict EU comparison).

A strategy $s$ is \emph{(weakly) dominant} for $i$ when it weakly
dominates every alternative; it is \emph{strictly dominant} when
the analogous statement holds for $\sPref$.  The library exposes
these as both EU-specific (\TT{IsDominant}, \TT{IsStrictDominant})
and preference-parameterized (\TT{IsDominantFor},
\TT{IsStrictDominantFor}).  The chain
\[
  \text{strict dominant} \;\Rightarrow\; \text{weak dominant} \;\Rightarrow\;
  \text{best response} \;\Rightarrow\; \text{Nash}
\]
is recorded modularly:\ each step is one \TT{simp}-able
proposition, with no hidden EU dependence.

\subsubsection{Iterated elimination and rationalizability}

Iterated elimination of strictly dominated strategies is encoded
as a sequence $\mathrm{Survives}_n$ of decreasing predicates: round
$0$ is true of every strategy, and round $n+1$ requires that the
strategy both survived round $n$ and is not strictly dominated,
restricted to round-$n$ survivors of all opponents.  A strategy is
\emph{rationalizable} when it survives every round.  The library
proves:
\begin{itemize}
\item Survival is monotone decreasing in rounds.
\item Every Nash strategy survives every round, hence is
rationalizable.
\item Every dominant strategy is rationalizable.
\item In a \emph{dominance-solvable} game, the unique surviving
profile is Nash.
\end{itemize}

\leancorr{Predicates and theorems in \TT{Concepts/Rationalizability.lean}
(\TT{Survives}, \TT{IsNash.survives}, \TT{IsNash.isRationalizable},
\TT{IsDominant.isRationalizable}); the dominance-solvable case in
\TT{Concepts/DominanceSolvable.lean}.}

\subsection{Correlated and coarse correlated equilibria}
\label{sec:solution-concepts:correlated}

Aumann's correlated equilibrium~\citep{Aumann1974} interprets
$\mu \in \PMF{\prod_i \Strat_i}$ as a public correlation device
that recommends a private action $s = \mu_i$ to each player.  The
device is incentive-compatible when no player gains in expectation
by ignoring the recommendation in favor of any function of it.
This is exactly $\mathcal{D}^{\mathrm{dep}}$-equilibrium under EU.
Coarse correlated equilibrium~\citep{HartMasColell2000} weakens the
deviation menu from arbitrary functions to constant strategies,
which yields $\mathcal{D}^{\mathrm{const}}$-equilibrium under EU.

The library proves the standard inclusions
\[
  \Dirac(\Pref\text{-Nash}) \;\subseteq\; \Pref\text{-CCE},
  \qquad
  \Pref\text{-CE} \;\subseteq\; \Pref\text{-CCE}.
\]
Here \TT{IsCorrelatedEqFor.toCoarseCorrelatedEqFor} records the CE
$\Rightarrow$ CCE inclusion; the Nash-to-CCE statement is the schema
instance $\mu = \Dirac_\Profile$.  In the EU-specific layer, existence of
correlated and coarse correlated equilibria for any finite
$\KG$ is the content of \S\ref{sec:theorems:correlated-existence}.

\subsection{Pareto efficiency, social welfare, and price of anarchy}
\label{sec:solution-concepts:welfare}

A profile $\Profile$ \emph{Pareto-dominates} $\tau$ for a pair
$(\Pref, \sPref)$ when every player weakly prefers $\Profile$'s
outcome distribution and at least one player strictly does;
$\Profile$ is \emph{Pareto-efficient} when no profile dominates
it.  The EU specialization uses $\Pref^{\EU}$ and
$\sPref^{\EU}$.  The library tracks the basic structural results
— irreflexivity, asymmetry — modulo the irreflexivity and
incompatibility hypotheses on $(\Pref, \sPref)$.

\emph{Social welfare} is the EU sum
$\Welfare(\KG, \Profile) = \sum_i \EU(\KG, \Profile, i)$ over a
finite player set.  The \emph{price of anarchy}
\citep{Koutsoupias1999, Roughgarden2005} compares the
worst-case Nash welfare to the optimum,
\begin{equation}
\label{eq:poa}
  \PoA(\KG) \;=\; \frac{\sup_\Profile \Welfare(\KG, \Profile)}
                       {\inf_{\Profile \in \Nash(\KG)} \Welfare(\KG, \Profile)},
\end{equation}
with the dual \emph{price of stability}
$\sup_\Profile \Welfare / \sup_{\Profile \in \Nash} \Welfare$.
The library formalizes the welfare quantities the ratios are built
from — optimal welfare, worst Nash welfare, best Nash welfare —
plus the comparison lemmas between them.  It does not yet package
the ratio-valued PoA/PoS definitions; those require explicit
positivity and non-empty-Nash hypotheses.

\leancorr{Welfare quantities in \TT{Concepts/PriceOfAnarchy.lean}
(\TT{optimalWelfare}, \TT{worstNashWelfare}, \TT{bestNashWelfare}).}

\subsection{Security strategies and minimax value}
\label{sec:solution-concepts:security}

Player $i$'s \emph{security level}~\citep{vNM1944} is the
guarantee value over a worst-case opponent profile,
\begin{equation}
\label{eq:security}
  \Sec(\KG, i) \;=\; \sup_{s \in \Strat_i}\;
    \inf_{\Profile_{-i}}\;
      \EU\bigl(\KG, \Profile_{-i}[i \mapsto s], i\bigr),
\end{equation}
and a \emph{security strategy} for $i$ achieves this supremum.
The library formalizes the worst-case-EU and security-level
quantities, proves that a security strategy exists in the finite
case, and proves that Nash payoffs are at least the security
level.  The two-player zero-sum theorem in
\S\ref{sec:theorems:minimax} is stated through a mixed Nash profile
with saddle-style guarantee inequalities.

\leancorr{Definitions and the IR/Nash bound in
\TT{Concepts/SecurityStrategy.lean} (\TT{worstCaseEU},
\TT{securityLevel}); individual-rationality interface in
\TT{Concepts/IndividualRationality.lean}.}

\subsection{Regret}
\label{sec:solution-concepts:regret}

Internal and external regret are recorded as positive parts of
EU differences.  External regret of $\Profile$ for player $i$ is
\begin{equation}
\label{eq:regret}
  \Regret^{\mathrm{ext}}(\KG, \Profile, i) \;=\;
    \max_{s' \in \Strat_i}\,
    \bigl(\EU(\KG, \Profile[i \mapsto s'], i) - \EU(\KG, \Profile, i)\bigr)^{+};
\end{equation}
internal regret replaces the constant $s'$ by a function of
$\Profile_i$ and is the dual of the $\mathcal{D}^{\mathrm{dep}}$
deviation family that defines correlated equilibrium.  The
library proves the structural fact that a profile distribution
$\mu$ has zero average external regret iff it is a CCE, and zero
average internal regret iff it is a CE — the standard
regret-equilibrium duality~\citep{HartMasColell2000}.

\subsection{Approximate equilibria}
\label{sec:solution-concepts:approx}

An \emph{$\varepsilon$-Nash equilibrium} relaxes the Nash
inequality by an additive $\varepsilon$ slack:\ every unilateral
deviation is at most $\varepsilon$-profitable.  Closure under
$\varepsilon$-relaxation is recorded for the deviation-family
schema, so the same definition applies to $\varepsilon$-CCE and
$\varepsilon$-CE.  We do not prove approximation theorems (e.g.,
LP-duality bounds for $\varepsilon$-CE in succinct games) in this
layer; we expose the definitions as the natural targets for
downstream complexity-theoretic work.

\leancorr{\TT{Concepts/ApproximateNash.lean}.}

\subsection{Potential games}
\label{sec:solution-concepts:potential}

A function $\Pot : \prod_i \Strat_i \to \RR$ is an \emph{exact
potential}~\citep{MonderShapley1996} when, for every player $i$
and unilateral deviation $\Profile \to \Profile[i \mapsto s']$,
\begin{equation}
\label{eq:potential}
  \EU(\KG, \Profile[i \mapsto s'], i) - \EU(\KG, \Profile, i)
  \;=\;
  \Pot(\Profile[i \mapsto s']) - \Pot(\Profile);
\end{equation}
it is an \emph{ordinal potential} when the equality is replaced
by an iff between the two sides being positive.  The library
proves:
\begin{itemize}\itemsep2pt
\item Every exact potential is an ordinal potential.
\item Every (local or global) maximizer of an ordinal potential
is Nash; hence finite ordinal potential games admit pure Nash
equilibria as a direct consequence (any maximizer of $\Pot$ on
the finite product $\prod_i \Strat_i$ is Nash).
\item Best-response dynamics terminates in finitely many steps in
a finite ordinal potential game.
\end{itemize}
Team games (\S\ref{sec:solution-concepts:team}) are an instance:\
they are exact potential games with $\Pot = u_{\mathrm{any}}$.

\leancorr{\TT{Concepts/PotentialGame.lean} (\TT{IsExactPotential},
\TT{IsOrdinalPotential}, \TT{IsExactPotential.toOrdinal},
\TT{IsOrdinalPotential.nash\_of\_maximizer});
\TT{Concepts/PotentialFIP.lean} and
\TT{Concepts/PotentialWellFounded.lean} for the dynamics;
\TT{Concepts/PotentialTeam.lean} for the team-game corollary
(\TT{IsTeamGame.isExactPotential}).}

\subsection{Zero-sum, constant-sum, team, and symmetric games}
\label{sec:solution-concepts:team}

A kernel game is \emph{zero-sum} when $\sum_i \util(\omega)(i) =
0$ for every $\omega$, and \emph{constant-sum} for an
$\Outcome$-independent constant.  The library proves that the
Nash equilibria of a two-player constant-sum game give each
player guarantee inequalities against opponent deviations.

A kernel game is a \emph{team game} when all players have
identical utilities.  Team games are exact potential games with
$\Pot = \util_{\mathrm{any}}$; consequently every welfare optimum
is a Nash equilibrium, and the price of stability is $1$.

A kernel game is \emph{symmetric} when permutations of the player
type act trivially on outcome distributions and utility (modulo
the same permutation).  The library proves the structural lemmas
that permuting a Nash profile by a symmetry yields a Nash profile,
and that all players obtain equal EU at a symmetric profile.
Existence of a \emph{symmetric} mixed Nash equilibrium (i.e.\ a
fixed point of the gain map within the symmetric subspace) is a
classical corollary of mixed Nash existence~\citep{Nash1951}; the
corollary is not packaged in the library and is listed in
\S\ref{sec:future} as a natural next step.  \emph{Evolutionary
stability}~\citep{MaynardSmithPrice1973} is formalized as the
appropriate strict-dominance condition on a symmetric two-player
game.

\leancorr{\TT{Concepts/ZeroSum.lean}, \TT{Concepts/ConstantSum.lean},
\TT{Concepts/ZeroSumNash.lean}, \TT{Concepts/ConstantSumNash.lean},
\TT{Concepts/TeamGame.lean}, \TT{Concepts/SymmetricGame.lean}
(\TT{isSymmetricEU\_nash\_perm}, \TT{isSymmetricEU\_eu\_eq}),
\TT{Concepts/EvolutionaryStability.lean}.}

\subsection{Common knowledge and individual rationality}
\label{sec:solution-concepts:misc}

Two further behavioral concepts are formalized for downstream
use.  An event is \emph{common knowledge} among a group of agents
when everyone knows it, everyone knows everyone knows it, and so
on for every finite depth of nested knowledge.  Aumann's
operator~\citep{Aumann1976} captures this fixed point on a finite
partition structure (each agent has an information partition over
the state space; common knowledge of an event $E$ at state $s$
holds iff $E$ contains the smallest set in the meet of all
partitions that contains $s$).  \emph{Individual rationality},
the second concept, is a participation constraint:\ a profile is
individually rational for player $i$ when $i$'s payoff is at least
their security level (\S\ref{sec:solution-concepts:security}).
Both are formal scaffolding for the mechanism-design layer
(\S\ref{sec:mechanism-auctions}).

\leancorr{\TT{Concepts/CommonKnowledge.lean};
\TT{Concepts/IndividualRationality.lean}.}
